\subsection{Bilinear maps.}

In this no., A and B denote two rings, E a left A-module, F a right B-module, and G an (A, B)-bimodule, that is, a commutative group equipped with a left A-module structure and a right B-module structure such that one has \((ag)b = a(gb)\) for all \(a \in A\) , \(b \in B\) , \(g \in G\) .

DEFINITION 1. --- One says that a map \(\Phi\) of the product \(\mathbf{E} \times \mathbf{F}\) in \(\mathbf{G}\) is bilinear if it satisfies the following conditions:

\begin{enumerate}
\def\labelenumi{(\arabic{enumi})}
\tightlist
\item
  \(\Phi(x + x', y) = \Phi(x, y) + \Phi(x', y)\)
\end{enumerate}

for all \(x \in \mathbf{E}\) , \(x' \in \mathbf{E}\) , \(y \in \mathbf{F}\) ;

\begin{enumerate}
\def\labelenumi{(\arabic{enumi})}
\setcounter{enumi}{1}
\tightlist
\item
  \(\Phi(x, y + y') = \Phi(x, y) + \Phi(x, y')\)
\end{enumerate}

for all \(x \in \mathbf{E}, y \in \mathbf{F}, y' \in \mathbf{F}\) ;

\begin{enumerate}
\def\labelenumi{(\arabic{enumi})}
\setcounter{enumi}{2}
\item
  \(\Phi(ax, y) = a\Phi(x, y)\) for all \(a \in \mathbf{A}\) , \(x \in \mathbf{E}\) , \(y \in \mathbf{F}\) ;
\item
  \(\Phi(x, yb) = \Phi(x, y)b\) for all \(x \in \mathbf{E}, y \in \mathbf{F}, b \in \mathbf{B}\) .
\end{enumerate}

The tensor product \(E \otimes_{Z} F\) is canonically endowed with a structure of (A, B)-bimodule characterized by \(a(x \otimes y)b = ax \otimes yb\) (Chap. III, 2nd ed., App. II, no. 3), and the data of a bilinear map \(\Phi\) of \(E \times F\) into G is equivalent to that of a map \(\Psi\) of \(E \otimes_{Z} F\) into G that is a homomorphism for the structures of (A, B)-bimodules and that satisfies \(\Psi(x \otimes y) = \Phi(x, y)\) for all \(x \in E\) and \(y \in F\) .

The conditions imposed on \(\Phi\) by Definition 1 mean that the partial maps \(d_{\Phi}(y): x \to \Phi(x, y)\) and \(s_{\Phi}(x): y \to \Phi(x, y)\) are respectively an A-linear map from E to G and a B-linear map from F to G. Let us equip the commutative group \(\mathcal{L}_{\mathbf{A}}(\mathbf{E}, \mathbf{G})\) (resp. \(\mathcal{L}_{\mathbf{B}}(\mathbf{F}, \mathbf{G})\) ) with the structure of a right B-module (resp. left A-module) defined by \(ub(x) = u(x). b\) ( \(u \in \mathcal{L}_{\mathbf{A}}(\mathbf{E}, \mathbf{G}), x \in \mathbf{E}, b \in \mathbf{B}\) ) (respectively \(av(y) = a. v(y)\) ( \(a \in \mathbf{A}, v \in \mathcal{L}_{\mathbf{B}}(\mathbf{F}, \mathbf{G}), y \in \mathbf{F}\) )). Then conditions (1) through (4) are respectively equivalent to:

\[
s _ {\Phi} (x + x ^ {\prime}) = s _ {\Phi} (x) + s _ {\Phi} (x ^ {\prime})\tag{1'}
\]

\[
d _ {\Phi} (y + y ^ {\prime}) = d _ {\Phi} (y) + d _ {\Phi} (y ^ {\prime})\tag{2'}
\]

\[
s _ {\Phi} (a x) = a. s _ {\Phi} (x)\tag{3'}
\]

\[
d _ {\Phi} (y b) = d _ {\Phi} (y). b,\tag{4'}
\]

for all \(x, x'\) in E, \(y, y'\) in F, \(a \in A\) , \(b \in B\) ; in other words, the map \(d_{\Phi}\) of F into \(\mathcal{L}_{A}(E, G)\) is B-linear, and the map \(s_{\Phi}\) from E into \(\mathcal{L}_{B}(F, G)\) is A-linear. We have, by definition,

\begin{enumerate}
\def\labelenumi{(\arabic{enumi})}
\setcounter{enumi}{4}
\tightlist
\item
  \(\Phi(x, y) = d_{\Phi}(y)(x) = s_{\Phi}(x)(y)\) for all \(x \in \mathbf{E}, y \in \mathbf{F}\) .
\end{enumerate}

DEFINITION 2. --- Given a bilinear map \(\Phi\) from E × F to G, the map \(d_{\Phi}\) of F into \(\mathfrak{L}_{\mathrm{A}}(\mathrm{E}, \mathrm{G})\) (resp. the map \(s_{\Phi}\) from E into \(\mathfrak{L}_{\mathrm{B}}(\mathrm{F}, \mathrm{G})\) ) characterized by (5) is called the linear map associated on the right (resp. on the left) to \(\Phi\) .

Conversely, the data of a B-linear map \(d\) of F into \(\mathfrak{L}_{\mathbf{A}}(\mathbf{E},\mathbf{G})\) (resp. of an A-linear map \(s\) from E into \(\mathfrak{L}_{\mathbf{B}}(\mathbf{F},\mathbf{G})\) ) determines uniquely, by the formula

\[
\Phi (x, y) = d (y) (x) \quad (\text { resp. } \Phi (x, y) = s (x) (y))
\]

a bilinear map \(\Phi\) of \(\mathbf{E} \times \mathbf{F}\) into \(\mathbf{G}\) , of which \(d\) (resp. \(s\) ) is the linear map associated on the right (resp. on the left).

DEFINITION 3. --- A bilinear map \(\Phi\) from E × F to G is said to be degenerate on the right (resp. on the left) if there exists a nonzero element \(y_0\) of F (resp. \(x_0\) of E) such that \(\Phi(x, y_0) = 0\) for all \(x \in \mathbf{E}\) , (resp. \(\Phi(x_0, y) = 0\) for all \(y \in \mathbf{F}\) ). One says that \(\Phi\) is degenerate if it is degenerate on the right or if it is degenerate on the left.

For \(\Phi\) to be nondegenerate on the right (resp. on the left) it is necessary and sufficient that the linear map associated on the right (resp. on the left) to \(\Phi\) be injective; to say that \(\Phi\) is nondegenerate therefore means that the associated linear maps \(d_{\Phi}\) and \(s_{\Phi}\) are both injective.

Let \((e_i)_{i\in I}\) and \((f_k)_{k\in K}\) be two families of elements of E and F, and let \((a_i)_{i\in I}\) and \((b_k)_{k\in K}\) be two families of elements of A and B, all but a finite number of which are zero. It follows from equalities (1) to (4), by induction on the number of nonzero coefficients, that one has

\[
\Phi (\sum_ {i} a _ {i} e _ {i}, \sum_ {k} f _ {k} b _ {k}) = \sum_ {i, k} a _ {i} \Phi (e _ {i}, f _ {k}) b _ {k}.\tag{6}
\]

If \((e_{i})\) and \((f_{k})\) are systems of generators of the modules E and F, \(\Phi\) is therefore completely determined by the elements \(g_{ik} = \Phi(e_{i}, f_{k})\) . If \((e_{i})\) and \((f_{k})\) are bases of E and F and elements \(g_{ik}\) of G \((i \in I, k \in K)\) are given, then the formula

\[
\Phi (\sum_ {i} a _ {i} e _ {i}, \sum_ {k} f _ {k} b _ {k}) = \sum_ {i, k} a _ {i} g _ {i k} b _ {k}\tag{6'}
\]

defines a map from E × F to G, which is bilinear and satisfies \(\Phi(e_{i}, f_{k}) = g_{ik}\) . When \((e_{i})\) and \((f_{k})\) are finite bases, we say that \((\Phi(e_{i}, f_{k}))\) is the matrix of \(\Phi\) with respect to these bases.

The bilinear maps from E × F into G obviously form a subgroup of the additive group of maps from E × F into G. On the other hand, let a (resp. b) be an element of the center of A (resp. B); then the map \(a\Phi b\) from E × F to G defined by \((a\Phi b)(x, y) = a \cdot \Phi(x, y) \cdot b\) is bilinear. The set of bilinear maps from E × F to G is thus endowed with a bimodule structure over the centers of A and B.

Let E' (resp. F') be a left A-module (resp. a right B-module), u (resp. v) a homomorphism from E to E' (resp. from F to F'), and \(\Phi'\) a bilinear map from \(E' \times F'\) into G. The inverse image of \(\Phi'\) (relative to u and v) is the bilinear map \(\Phi\) from E × F to G defined by

\[
\Phi (x, y) = \Phi^ {\prime} (u (x), v (y)) \quad (x \in \mathrm{E}, y \in \mathrm{F}).
\]

One verifies easily that one has

\[
d _ {\Phi} (y) = d _ {\Phi^ {\prime}} (\nu (y)) \circ u \quad \text { and } \quad s _ {\Phi} (x) = s _ {\Phi^ {\prime}} (u (x)) \circ \nu
\]

for all \(x \in \mathbf{E}, y \in \mathbf{F}\) .

Let \(\Phi\) be a bilinear map from \(\mathbf{E} \times \mathbf{F}\) into \(\mathbf{G}\) , and \(h\) a homomorphism (for the structures of (A, B)-bimodules) of \(\mathbf{G}\) into another (A, B)-bimodule \(\mathbf{G}'\) . Then \(h \circ \Phi\) is a bilinear map from \(\mathbf{E} \times \mathbf{F}\) into \(\mathbf{G}'\) .

